\documentclass[10pt,a4paper]{amsart}
\usepackage[margin=19mm]{geometry}
\usepackage{amsmath,amssymb,mathtools}
\usepackage[scr=rsfs,cal=euler]{mathalfa}
\usepackage{xfrac}
\usepackage[unicode,hidelinks,bookmarks=false]{hyperref}
\newcommand{\ie}{i.e.\ }
\newcommand{\cf}{cf.\ }
\newcommand{\resp}{respectively\ }
\newcommand{\Subsection}{\subsection}
\providecommand{\nobreakdash}{\nobreak\hbox{-}}
\setlength{\emergencystretch}{2em}
\multlinegap0pt
\usepackage{bm}
\usepackage{bbm}
\usepackage{dsfont}
\usepackage{xspace}
\usepackage{cancel}
\usepackage{mathtools}
\usepackage{amsthm}
\makeatletter
\@ifundefined{theorem}{\newtheorem{theorem}{Theorem}}{}
\makeatother
\title[Constructions of Compact Dupin Hypersurfaces with Non-consta]{Constructions of Compact Dupin Hypersurfaces with Non-constant Lie Curvatures}
\author{Thomas E. Cecil}
\date{Source: arXiv:2509.21235v2 (CC BY 4.0)}
\begin{document}
\maketitle

\noindent\emph{Source and licence.} Constructions of Compact Dupin Hypersurfaces with Non-constant Lie Curvatures. Version arXiv:2509.21235v2 (CC BY 4.0), distributed under the Creative Commons Attribution 4.0 International licence, as recorded in the arXiv licence metadata for this exact version and shown on the abstract page https://arxiv.org/abs/2509.21235v2. Full rights notice: licenses/arxiv-2509.21235-CC-BY-4.0.txt.\par\smallskip

\noindent\textit{Editorial scope.} Counted unit: the Theorem whose source label is \texttt{thm:3.7.1} in the article (source0.tex lines 852--855, source label \texttt{thm:3.7.1}), delivered together with the article's own complete original proof of it (source0.tex lines 856--950, one \texttt{proof} environment of 95 lines). No mathematics is written, repaired or re-derived: statement and proof are the article's own text line for line. Required context retained uncounted: eq:3.7.13 (lines 826--829). Every definition, display and label of those lines is retained unchanged. Omitted from this package and not redistributed: 1--825, 830--851, 951--1261. Nothing inside a retained excerpt is omitted or altered: every retained body is delivered exactly as the article's lines are written, so no omission receipt of any kind accompanies this package. Citations: the retained excerpts carry no citation command at all, so no bibliography entry of the article is transcribed and no reference is redistributed. Reference adaptation: none was needed and none was made; every reference inside the delivered unit resolves inside it, so no unresolved reference or citation is produced. Printed slips kept literal: no printed word, symbol, hypothesis or number of the counted statement or of its proof was altered. Layout and class: the article's own class is \texttt{article} with packages including amsmath, amsthm, amssymb, euscript, verbatim, makeidx, graphicx, multicol, footmisc, latexsym, epic, eepic, pstricks; the wrapper is the lane's pinned \texttt{amsart} editorial wrapper at 10pt on A4, which restates the theorem-like environments the retained text needs and adds the article's own macro definitions that the retained text needs (none), with extra wrapper packages (bm, bbm, dsfont, xspace, cancel, mathtools). The delivered numbering is the unit's own. Assets: no retained span contains a figure environment, a graphics-inclusion command, a PDF-page inclusion, a table or a TikZ picture; no image file is copied into this package, the package declares no asset dependency and no image file is redistributed with it. Licence: version arXiv:2509.21235v2 of this article is distributed under the Creative Commons Attribution 4.0 International licence, as recorded in the arXiv licence metadata for this exact version and shown on the abstract page https://arxiv.org/abs/2509.21235v2. A full rights notice is delivered at licenses/arxiv-2509.21235-CC-BY-4.0.txt. Characters: every retained excerpt is pure ASCII, so the delivered engine, XeLaTeX with its default Latin Modern text font, reports no missing character.
\begin{equation} 
\label{eq:3.7.13}
u = \frac{wz}{\sqrt{2(1-t)}}, \quad v = \sqrt{(1-t)/2}\  z, \quad z \in S^3.
\end{equation}

\begin{theorem}
\label{thm:3.7.1} 
Let $M$ be a compact, connected submanifold of $S^4$.  If $M$ is taut in $S^4$, then $h^{-1}(M)$ is taut in $S^7$.
\end{theorem}

\begin{proof}
Since both $M$ and $h^{-1}(M)$ lie in spheres, tautness of $h^{-1}(M)$ in $S^7$ is equivalent to tightness of
$h^{-1}(M)$ in ${\bf R}^8$, i.e., every nondegenerate linear 
height function in ${\bf R}^8$ has the minimum number
of critical points.  We write linear height functions in ${\bf R}^8$ in the form
\begin{equation} 
\label{eq:3.7.15}
f_{ab} (u,v) = \Re (au + bv) = (\bar{a}, \bar{b}) \cdot (u,v), \quad (a,b) \in S^7.
\end{equation}
This is the height function in the direction $(\bar{a}, \bar{b})$.  We want to determine when the point $(u,v)$
is a critical point of $f_{ab}$.  Without loss of generality, we may assume that $(u,v)$ lies in a local
trivialization of the form \eqref{eq:3.7.13} when making local calculations.  Let $x = (w,t)$ be a point of
$M \subset S^4$, and let 
\begin{displaymath}
(x,z) = (w,t,z)
\end{displaymath}
be a point in the fiber $h^{-1}(x)$.  The tangent space to $h^{-1}(M)$
at $(x,z)$ can be decomposed as $T_x M \times T_z S^3$.  We first locate the critical points of the restriction of 
$f_{ab}$ to the fiber through $(x,z)$.  By equations \eqref{eq:3.7.13} and \eqref{eq:3.7.15}, we have
\begin{eqnarray} 
\label{eq:3.7.16}
f_{ab} (w,t,z) & = & \Re \left(\frac{awz}{\sqrt{2(1-t)}} + bz \sqrt{(1-t)/2} \right) \\
& = & \Re (\alpha (w,t) z) = \alpha (w,t) \cdot \bar{z}, \nonumber
\end{eqnarray}
where
\begin{displaymath}
\alpha (w,t) = \frac{aw}{\sqrt{2(1-t)}} + b \sqrt{(1-t)/2}.
\end{displaymath}
This defines the map $\alpha$ from $S^4$ to ${\bf H}$.  If $Z$ is any tangent vector to $S^3$ at $z$, we write
$Zf_{ab}$ for the derivative of $f_{ab}$ in the direction $(0,Z)$.  Then
\begin{equation} 
\label{eq:3.7.17}
Zf_{ab} = \alpha (w,t) \cdot \bar{Z}
\end{equation}
at $(x,z)$.  Now there are two cases to consider.  First, if $\alpha (w,t) \neq 0$, then in order to have 
$Zf_{ab} = 0$ for all $Z \in T_z S^3$, we must have
\begin{equation} 
\label{eq:3.7.18}
\bar{z} = \pm \frac{\alpha (w,t)}{|\alpha (w,t)|}.
\end{equation}
So the restriction of $f_{ab}$ to the fiber has exactly two critical points with corresponding values
\begin{equation} 
\label{eq:3.7.19}
\pm |\alpha (w,t)|.
\end{equation}
The second case is when $\alpha (w,t) = 0$.  Then the restriction of $f_{ab}$ to the fiber is identically zero
by equation \eqref{eq:3.7.16}.  In both cases the function,
\begin{displaymath}
g_{ab} (w,t) = |\alpha (w,t)|^2,
\end{displaymath}
satisfies the equation
\begin{displaymath}
g_{ab} (w,t) = f_{ab}^2 (w,t,z),
\end{displaymath}
at the critical point.  The key in relating this fact to information about the submanifold $M$ is to note that
\begin{eqnarray} 
\label{eq:3.7.20}
g_{ab} (w,t) & = & |\alpha (w,t)|^2 = \frac{1}{2} \Re \{2 a \bar{b} w + (|a|^2 - |b|^2)t\} + \frac{1}{2}(|a|^2 + |b|^2)
\nonumber \\
 & = & \frac{1}{2} + \frac{1}{2} ((w,t) \cdot (2 \bar{a} b, |a|^2 - |b|^2)) \\
& = & \frac{1}{2} + \frac{1}{2} \ell_{ab} (w,t), \nonumber
\end{eqnarray}
where $\ell_{ab}$ is the linear height function on ${\bf R}^5$ in the direction
\begin{displaymath}
(2\bar{a} b, |a|^2 - |b|^2) = h (\bar{a}, \bar{b}).
\end{displaymath}
This shows that $g_{ab} (w,t) = 0$ if and only if $(w,t) = - h(\bar{a}, \bar{b})$.  Thus, if $- h(\bar{a}, \bar{b})$
is not in $M$, the restriction of $f_{ab}$ to each fiber has exactly two critical points of the form $(x,z)$, with $z$
as in equation \eqref{eq:3.7.18}.  For $X \in T_x M$, we write $Xf_{ab}$ for the derivative of $f_{ab}$ in the direction
$(X,0)$.  At the two critical points, we have
\begin{equation} 
\label{eq:3.7.21}
Xf_{ab} = d\alpha (X) \cdot \bar{z},
\end{equation}
\begin{equation} 
\label{eq:3.7.22}
Xg_{ab} = 2 d\alpha (X) \cdot \alpha (x) = \pm 2|\alpha (x)| (d\alpha (X)\cdot \bar{z}) = \pm 2|\alpha (X)| Xf_{ab}.
\end{equation}
Thus $(x,z)$ is a critical point of $f_{ab}$ if and only if $x$ is a critical point of $g_{ab}$.  By equation
\eqref{eq:3.7.20}, this happens precisely when $x$ is a critical point of $\ell_{ab}$.  We conclude that if 
$- h(\bar{a}, \bar{b})$ is not in $M$, then $f_{ab}$ has two critical points for every critical point of 
$\ell_{ab}$ on $M$.  The set of points $(a,b)$ in $S^7$ such that $- h(\bar{a}, \bar{b})$ belongs to $M$ has measure
zero.  If $(a,b)$ is not in this set, then $f_{ab}$ has twice as many critical points as the height function
$\ell_{ab}$ on $M$.  Since $M$ is taut, every nondegenerate 
height function $\ell_{ab}$ has $\beta (M;{\bf Z}_2)$ critical
points on $M$, where $\beta (M;{\bf Z}_2)$ is the sum of the 
${\bf Z}_2$-Betti numbers of $M$.  
Thus, except for a set of measure zero, every height function $f_{ab}$
has $2 \beta (M;{\bf Z}_2)$ critical points on $h^{-1}(M)$.  Since $h^{-1}(M)$ is diffeomorphic to $M \times S^3$,
we have 
\begin{displaymath}
\beta (h^{-1}(M);{\bf Z}_2) = \beta (M \times S^3;{\bf Z}_2) = 2 \beta (M;{\bf Z}_2).
\end{displaymath}
Thus, $h^{-1}(M)$ is taut in $S^7$. 
\end{proof}

\end{document}
